\section{系统测试与结果分析}

第三章给出了伏羲（Fuxi）的设计与实现细节，本章在该实现基础上对系统进行多维度测试。测试目标分为正确性、实时性、吞吐扩展性、参数敏感性以及通讯层压力承受能力五大类，覆盖从底层事件总线到顶层任务编排的完整链路。所有定量实验在同一台机器完成，配置：Apple M4 Pro（10 P 核 + 4 E 核 / 14 物理核）、48 GB 统一内存、macOS 26.0.1、Rust 1.93.1 stable、SQLite 3.51.0、Tokio 多线程运行时显式 \texttt{worker\_threads(8)}。吞吐采用 5-run median 抑制冷启动，延迟采用 500 独立样本统计 p50/p99，事件总线压测每 cell 5 秒稳态。完整原始数据见仓库 \texttt{deliverables/thesis-v2/benchmarks/v2-2026-05-07.md}。

\subsection{测试目标与方法学}

正确性通过 Rust 单元测试与跨 crate 集成测试验证；实时性通过任务派发与事件总线广播的端到端延迟测量验证，对应分布式系统因果序传递的经典指标\cite{lamport1978time}；吞吐扩展性通过固定任务粒度、扫 Worker 数量的水平扩展实验验证，思路与 MapReduce\cite{dean2004mapreduce} 一致；参数敏感性通过固定其他变量、扫描 \texttt{poll\_ms} 的消融实验验证；通讯层压力承受能力通过事件总线纯压测验证，参考了 Kafka 等持久化消息系统的纯通讯极限测试\cite{kreps2011kafka}。性能测试以独立 \texttt{cargo bench} 二进制 release profile 编译，所有数据由 bench 写入 markdown 报告，论文中所有图表由 matplotlib 脚本基于该报告自动生成，规避「论文里数字与仓库里数字不一致」常见问题。

测量方法学方面采用以下控制手段：吞吐与扩展性实验「5 次重复取中位数」抑制冷启动，对应顶刊 systems 论文常见做法；延迟实验在采集 500 样本前先丢弃 50 个 warmup；测试期间关闭 Spotlight 索引、Time Machine、IDE 索引等可能消耗 CPU 的进程并选夜间安静时段；Tokio 多线程运行时显式 \texttt{worker\_threads(8)} 固定避免不同实验线程数变化导致结果不可比；时钟使用 \texttt{std::time::Instant}（macOS 上为 \texttt{mach\_absolute\_time}，分辨率约 1 ns），但延迟统计经 \texttt{SystemTime} 的 ns→us 整数除法截断，凡 p50 报告为 0.0 us 处实际含义为「< 1 us」。基于 500 个独立样本，p99 估计的 95\% Wilson 置信区间约 ±2\%，远小于本章关注的数量级差异。

baseline 选择上，本章未引入外部系统作为实测对照——公开的 AI Agent 编排系统多以 Python 实现，跨语言、跨抽象层级的 baseline 容易让数字失去工程意义（差异主要来自语言运行时与同步异步模型，而非编排层设计取舍），且容易让读者形成「Rust 比 Python 快」的浅层结论。本章采用\textbf{理论上限对照}（每 cell 直接给出 \texttt{tasks\_n × 1000 / job\_sleep\_ms} 作为零开销完美调度上限），使读者可直接读出 Fuxi 与这一理想上限之间的工程损耗。

\subsection{测试用例覆盖}

Fuxi 的集成测试在每一个 crate 中按其对外暴露的能力组织，完整覆盖如表~\ref{tab:test-coverage}。总计 27 个集成测试文件分布在 13 个 crate，覆盖通讯、编排、执行、观测与支撑模块的全部主要路径。

\begin{table}[H]
\centering
\caption{系统测试覆盖情况}
\label{tab:test-coverage}
\footnotesize
\begin{tabularx}{\textwidth}{p{2.7cm} p{6.0cm} X}
\toprule
模块 & 主要测试文件 & 测试目标 \\
\midrule
fuxi-a2a & \texttt{roundtrip.rs} & A2A wire 编解码与往返一致性 \\
fuxi-events & \texttt{integration.rs} & 事件发布、订阅与回放 \\
fuxi-agent-cc & \texttt{fixture\_stream.rs}、\texttt{real\_cc\_smoke.rs} & Claude Code 输出解析与真实启动 \\
fuxi-agent-codex & \texttt{fixture\_stream.rs}、\texttt{real\_codex\_smoke.rs} & Codex 输出解析与真实启动 \\
fuxi-cli & \texttt{chaos\_resilience.rs} 等 5 文件 & CLI、网关重启、分布式与召贤 e2e \\
fuxi-im & \texttt{router\_smoke.rs} 等 3 文件 & IM API、WebSocket 与认证 \\
fuxi-firehose & \texttt{hub\_roundtrip.rs} & Firehose 事件 hub 往返 \\
fuxi-memory & \texttt{oracle\_roundtrip.rs} 等 4 文件 & 长期记忆、可复用模式与抽取 \\
fuxi-orchestrator & \texttt{dispatch.rs}、\texttt{deliverable\_handoff.rs} & 任务派发与交付物交接 \\
fuxi-scheduler & \texttt{e2e\_scheduler.rs} & 定时触发链路 \\
fuxi-skills & \texttt{loader.rs} 等 4 文件 & 角色加载、暂存、登记与模板 \\
fuxi-workspace & \texttt{integration.rs} & 工作区创建与隔离 \\
\bottomrule
\end{tabularx}
\end{table}

集成测试覆盖广度并不等价于行覆盖率指标——Rust 异步生态下 \texttt{cargo-llvm-cov} 对 \texttt{tokio::select!}、\texttt{tokio::spawn} 内部分支的归因尚不完整，强行追求百分数易诱导写「为覆盖而覆盖」的低价值测试。本章选择以「关键路径必有 e2e」「关键 wire 必有 round-trip」「关键 FSM 必有状态迁移枚举」三项硬规则替代覆盖率门槛。最近一次提交中 \texttt{cargo test --workspace --all-targets} 完整通过，IM PWA 前端 \texttt{pnpm test}、typecheck 与 lint 也已通过。集成测试之外还使用 property-based testing（\texttt{proptest} 对 wire round-trip 不变量验证）、回归 fixture 沉淀（每次发现新 bug 优先沉淀一条 fixture）与 smoke 测试（\texttt{smoke-p0.sh} 汇总系统启动—对话—门客派发—回收的完整链路）三类辅助验证手段。

\subsection{吞吐量、扩展性与端到端延迟分解}

吞吐量测试基于第二章 SLA 节给出的端到端延迟分解 $L_{e2e} = L_{enq} + L_{disp} + L_{exec} + L_{evt}$ 与吞吐损耗 $\eta = 1 - T/T_{ideal}$。本节把通讯与调度相关的三个分量合记 $T_{comm} \equiv L_{enq} + L_{disp} + L_{evt}$、把任务执行单独写作 $T_{exec} \equiv L_{exec}$。基线吞吐结果如表~\ref{tab:throughput}。

\begin{table}[H]
\centering
\caption{分布式任务吞吐量基线测试结果（5-run median）}
\label{tab:throughput}
\begin{tabular}{rrrrrrr}
\toprule
worker\_n & job\_sleep & tasks\_n & median\_wall\_ms & tasks/s & 理论上限 & Fuxi 损耗 \\
\midrule
1 & 10 ms & 100 & 1227 & 81.50 & 100.00 & 18.5\% \\
1 & 100 ms & 50 & 5123 & 9.76 & 10.00 & 2.4\% \\
1 & 1000 ms & 10 & 10040 & 1.00 & 1.00 & 0.4\% \\
4 & 10 ms & 400 & 1212 & 330.03 & 400.00 & 17.5\% \\
4 & 100 ms & 200 & 5110 & 39.14 & 40.00 & 2.2\% \\
4 & 1000 ms & 40 & 10033 & 3.99 & 4.00 & 0.3\% \\
8 & 10 ms & 800 & 1203 & 665.00 & 800.00 & 16.9\% \\
8 & 100 ms & 400 & 5124 & 78.06 & 80.00 & 2.4\% \\
8 & 1000 ms & 80 & 10046 & 7.96 & 8.00 & 0.5\% \\
\bottomrule
\end{tabular}
\end{table}

表中三组层次分明的现象：10 ms 任务下损耗稳定在 17\%--19\%，100 ms 损耗骤降到 2.2\%--2.4\%，1000 ms 损耗降至 0.3\%--0.5\%。三段递减的物理原因是 $T_{comm}$（通讯与调度开销）在不同任务粒度下的占比差异：$T_{comm}$ 包含 axum HTTP 往返（约 4--6 ms）、Tokio 任务调度（约 2--3 ms）、SQLite 状态写入（约 1--2 ms）以及 Worker 拉取轮询（poll\_ms 平均一半）。具体地，10 ms 任务下 $T_{comm}/T_{e2e} \approx 55\%$，但因多个任务在同一时段并发流水化（dispatcher 在第一个任务执行时同时为下一个任务做派发），最终 wall-clock 损耗约 18\%；100 ms 下 $T_{comm}/T_{e2e} \approx 11\%$，1000 ms 下仅 1.2\%。

\textbf{1 worker 与 8 worker 损耗几乎相同}（10 ms 任务下分别为 18.5\% 与 16.9\%），这一现象在传统串行系统中几乎不可能出现——通常 worker 数增加会加剧 dispatcher 串行化压力让损耗单调上升。Fuxi 损耗保持平稳的原因有二：dispatcher 自身用 Tokio 异步路径处理请求，单 dispatcher 在本机能稳定承担 1000+ TPS 远超实际负载；事件总线 fan-out 是 O(1) 的（每个订阅者只递增引用计数），即使 worker 数增加每个 worker 收到任务通知的延迟也几乎不变。这两点合起来让 Fuxi 在小规模水平扩展下表现出「损耗常数化」的优良性质。

为进一步验证水平扩展能力，固定 \texttt{job\_sleep=10 ms}、\texttt{poll\_ms=5}、每 worker 跑 100 task，扫描 worker 数 \{1, 2, 4, 8, 16\}，结果如表~\ref{tab:scalability} 与图~\ref{fig:scalability}。

\begin{table}[H]
\centering
\caption{Worker 扩展性测试结果（5-run median）}
\label{tab:scalability}
\begin{tabular}{rrrrrr}
\toprule
worker\_n & tasks\_n & median\_wall\_ms & tasks/s & 理论上限 & 扩展效率 \\
\midrule
1 & 100 & 1227 & 81.50 & 100.00 & 81.5\% \\
2 & 200 & 1230 & 162.60 & 200.00 & 81.3\% \\
4 & 400 & 1212 & 330.03 & 400.00 & 82.5\% \\
8 & 800 & 1203 & 665.00 & 800.00 & 83.1\% \\
16 & 1600 & 1242 & 1288.24 & 1600.00 & 80.5\% \\
\bottomrule
\end{tabular}
\end{table}

\begin{figure}[H]
\centering
\includegraphics[width=0.95\linewidth]{fig-5-1-scalability.pdf}
\caption{Worker 数对吞吐量（左）与扩展效率（右）的影响}
\label{fig:scalability}
\end{figure}

需要先解释 wall-clock 几乎不随 worker 数下降的现象：本节固定「每 worker 100 task」，1 worker 与 16 worker 场景下每个 worker 各自仍要串行跑完 100 task，关键路径长度由单 worker 自己的 100 task 串行执行决定而非由全部 1600 task 决定。换言之，本节衡量的是「dispatcher + 事件总线在 1$\to$16 worker 范围内能否保持 per-worker 关键路径不被自身串行化拖累」——若 dispatcher 在 16 worker 下成为瓶颈 per-worker 关键路径会被拉长。实测 wall-clock 几乎不变说明在被测机上 dispatcher 在 16 worker × 10 ms 任务粒度下尚未饱和。扩展效率 80.5\%--83.1\% 的物理来源结构化拆解为三层：P 核与 E 核的差异（worker 数超过 P 核数 10 时新增 worker 部分被调度到 E 核，单核吞吐约为 P 核 50\%--70\%）、L3 cache 与内存带宽竞争（多 worker 共享同一片 L3 触发 L2/L3 miss）、dispatcher 单点的串行化代价（实测 16 worker 下约 35\% CPU，未成瓶颈）。上述分量比例为定性归因，未做 profiler 独立验证。

把吞吐与延迟实验按 SLA 公式分量分解可得到不同任务粒度下端到端延迟相对构成（图~\ref{fig:e2e-breakdown}）。需明确：图中 $L_{enq}$ 与 $L_{disp}$ 在 Fuxi 当前实现里共享同一次 SQLite 事务难以独立测量，按经验比例 3:7 对 task\_dispatch p50 = 32.93 ms 估算拆分。

\begin{figure}[H]
\centering
\includegraphics[width=0.95\linewidth]{fig-5-6-e2e-breakdown.pdf}
\caption{端到端延迟分解（10/100/1000 ms 任务粒度）}
\label{fig:e2e-breakdown}
\end{figure}

\textbf{三种工作负载下分量占比}：10 ms 任务下通讯与调度开销（约 33 ms）已超过执行开销，通讯占比 77\%；100 ms 任务下通讯占比降至 25\%；1000 ms 任务下通讯占比仅 3.2\%。当任务粒度大于 ~300 ms 时通讯开销在端到端延迟中占比降至 10\% 以下，对终端用户体验已不可感知。考虑到真实 AI Agent 场景中 LLM 推理 typical 1--10 秒、工具调用 typical 0.5--5 秒、用户对话回合典型 5--30 秒——单任务粒度均远超 300 ms 量级，Fuxi 当前的通讯开销在实际工作负载中不构成瓶颈。$L_{evt} = 0.08$ ms 在三种任务粒度下的占比都极小，目的不是优化端到端用户感知延迟（用户感知由 $L_{exec}$ 主导），而是为了让事件流在所有任务粒度下都「看起来是实时的」——若 event\_flow 与 task\_dispatch 同量级，观察端在显示进度时需要做 lookahead 缓冲与时间戳重排，UI 层复杂度显著上升。

\subsection{延迟分布与尾部分析}

延迟测试针对系统的两条关键通讯路径：任务派发链路（含 HTTP 往返、HMAC 验证、SQLite 写入、worker 选择与反序列化）与事件总线广播（同进程内 Tokio broadcast publish→subscribe 的 fan-out 路径，不经过网络/文件 I/O/序列化）。结果如表~\ref{tab:latency} 与图~\ref{fig:dispatch-latency}~\ref{fig:event-flow-latency}。

\begin{table}[H]
\centering
\caption{任务派发与事件流延迟测试结果（n=500）}
\label{tab:latency}
\begin{tabular}{lrrrrr}
\toprule
metric & sample\_n & min\_ms & p50\_ms & p99\_ms & max\_ms \\
\midrule
task\_dispatch & 500 & 11.64 & 32.93 & 36.66 & 37.43 \\
event\_flow & 500 & 0.04 & 0.08 & 0.26 & 0.66 \\
\bottomrule
\end{tabular}
\end{table}

\begin{figure}[H]
\centering
\includegraphics[width=0.7\linewidth]{fig-5-4-dispatch-latency.pdf}
\caption{任务派发延迟经验累积分布（n=500）}
\label{fig:dispatch-latency}
\end{figure}

\begin{figure}[H]
\centering
\includegraphics[width=0.7\linewidth]{fig-5-5-event-flow-latency.pdf}
\caption{事件总线广播延迟经验累积分布（n=500）}
\label{fig:event-flow-latency}
\end{figure}

任务派发 p99/p50 比值为 1.11，说明派发链路尾部抖动可控；事件总线广播 p99/p50 比值为 3.25，绝对差仅 0.18 ms（在亚毫秒量级 OS 时钟分辨率即足以让相对比值放大）。对 task\_dispatch p50 = 32.93 ms 做工程估算分解（基于第三章实现细节，未经 profiler 独立验证）：HTTP 往返约 10 ms，HMAC-SHA256 约 1 ms，Tokio 任务调度约 5 ms，SQLite 状态写约 6 ms，dispatcher 内部任务匹配约 4 ms，剩余 7 ms 来自 worker 接收事件、反序列化任务体并 ack 回 dispatcher。event\_flow p99 = 0.26 ms 的物理下界来自单次 Arc clone + 原子计数器递增（约 50 ns）、跨线程 future 唤醒（约 5--20 us）以及 broadcast 内部锁的获取释放（约 1 us 量级），p99 与下限 210 us 的差距可归因于 Tokio runtime 在 8 worker thread 之间负载均衡偶发的 task migration\cite{tokio2024broadcast}。

\textbf{event\_flow 与 task\_dispatch 三个数量级差异的工程意义}：事件流 p50 = 0.08 ms 与任务派发 p50 = 32.93 ms 相差约 410 倍。该差异反映了 Fuxi 在系统层级上的设计意图：通讯层应具备远低于业务层的延迟下界，使观测端能在任务执行同时近实时感知所有过程事件。如果事件流与任务派发同量级，Firehose、TUI、PWA 等观察端在显示进度时会出现明显滞后——这正是公理「真实时不轮询」试图避免的体验缺陷。task\_dispatch p99 = 36.66 ms 与 max = 37.43 ms 之间相差不到 1 ms，说明 500 样本中超过 p99 的样本紧密分布未出现孤立长尾；前序实验在另一份相同配置 run 中曾出现单个 227 ms 极端长尾样本归因于 macOS 后台调度突发（Spotlight 重索引、Time Machine 触发等）。即便保守按 200 ms 量级估计，单次任务派发的最坏情形相对一次完整玄女—门客交互（端到端用户感知 1--3 秒含 LLM 推理）占比仍 < 23\%，且 LLM 推理本身方差远大于此，因此 dispatch 长尾不会被用户感知。

\subsection{poll\_ms 参数消融}

为评估 Worker 拉取间隔 \texttt{poll\_ms} 对系统吞吐影响，固定 \texttt{worker\_n=4}、\texttt{tasks\_n=400}、\texttt{job\_sleep=10 ms}，扫描 \texttt{poll\_ms} $\in$ \{5, 10, 25, 50, 100\}，结果如表~\ref{tab:poll-scan} 与图~\ref{fig:poll-scan}。

\begin{table}[H]
\centering
\caption{\texttt{poll\_ms} 参数消融实验结果}
\label{tab:poll-scan}
\begin{tabular}{rrrrr}
\toprule
poll\_ms & median\_wall\_ms & tasks/s & 理论上限 & Fuxi 损耗 \\
\midrule
5 & 1285 & 311.28 & 400.00 & 22.2\% \\
10 & 1299 & 307.93 & 400.00 & 23.0\% \\
25 & 1242 & 322.06 & 400.00 & 19.5\% \\
50 & 1225 & 326.53 & 400.00 & 18.4\% \\
100 & 1235 & 323.89 & 400.00 & 19.0\% \\
\bottomrule
\end{tabular}
\end{table}

\begin{figure}[H]
\centering
\includegraphics[width=0.7\linewidth]{fig-5-2-poll-scan.pdf}
\caption{\texttt{poll\_ms} 参数对吞吐与损耗的影响}
\label{fig:poll-scan}
\end{figure}

\texttt{poll\_ms} 从 5 ms 到 100 ms 区间内吞吐与损耗差异不超过 3.6 个百分点。这一现象在 sustained 与 burst 两类场景下应做差异化解读：sustained 负载下 Worker 几乎不进入 idle 等待，每完成一个任务即可在下一个心跳前领到下一个任务，\texttt{poll\_ms} 几乎不影响关键路径；burst 负载下 \texttt{poll\_ms} 决定了从 idle 唤醒到拿到任务的期望延迟（期望值为 \texttt{poll\_ms}/2）。本节实验属于前者——400 个 10 ms 任务集中下发全程几乎无 idle，因此 \texttt{poll\_ms} 的 20 倍变化只在吞吐上带来 3.6 个百分点的差异。\texttt{poll\_ms=5} 反而损耗更高（22.2\%）的成因是过频 polling 在 sustained 负载下的副作用：每次 poll 都是一次到 dispatcher 的 HTTP 询问即使返回「无新任务」也消耗请求处理槽位与 worker CPU 上下文切换。综合 sustained 低 CPU 开销与 burst 响应性，本文推荐 Fuxi 默认 \texttt{poll\_ms} 取 25--50 ms。

需说明 Fuxi 内部「事件驱动 vs 轮询」的层次区分：dispatch 路径采用 worker 主动 poll（间隔由 \texttt{poll\_ms} 控制），observation 路径走 EventBus 订阅（无 poll）。本节扫描验证的是前者在 sustained throughput 场景下采用 poll-only 设计的性能代价；event\_flow p99 = 0.26 ms 显示事件总线纯通知路径已可在亚毫秒级完成扇出。当前 dispatch 路径采用 poll-only 模型在 sustained 场景下未表现为主要瓶颈，但本文未实测 push-based dispatch 对照版本——「事件驱动派发能否降低 burst 场景延迟」属于后续工作。当前选择保留 poll-only dispatch 的理由是工程权衡：sustained throughput 是日用平台主导工作负载，代码与故障模型也比双通道（push + 兜底）简单；为压低 burst 单任务派发延迟引入 push 机制需在跨节点心跳与网络分区下处理通知丢失复杂度，与本工作「单机优先」定位不匹配。

\subsection{事件总线压力承受能力}

事件总线是系统通讯核心组件。本节设计 4 sub × 3 rate × 2 payload = 24 cell 的事件总线纯压测，每 cell 5 秒稳态。关键 cell 如表~\ref{tab:bus-stress} 与图~\ref{fig:bus-stress}。

\begin{table}[H]
\centering
\caption{事件总线纯压测关键 cell（每 cell 5 秒稳态）}
\label{tab:bus-stress}
\begin{tabular}{rrlrrrr}
\toprule
sub\_n & rate(ev/s) & payload & publish\_tps & p50\_us & p99\_us & drops \\
\midrule
1 & 1k & small & 1000 & < 1 & < 1 & 0 \\
4 & 1k & small & 1000 & < 1 & < 1 & 0 \\
16 & 1k & small & 1000 & < 1 & < 1 & 0 \\
64 & 1k & small & 1000 & < 1 & < 1 & 0 \\
16 & 100k & small & 100000 & < 1 & < 1 & 0 \\
16 & 100k & large & 100000 & < 1 & < 1 & 0 \\
64 & 10k & small & 10000 & < 1 & < 1 & 0 \\
64 & 10k & large & 10000 & < 1 & 0.2 & 0 \\
64 & 100k & small & 100000 & 163.5 & 278.4 & 13599792 \\
64 & 100k & large & 100000 & 163.5 & 278.7 & 13598222 \\
\bottomrule
\end{tabular}
\end{table}

\begin{figure}[H]
\centering
\includegraphics[width=0.85\linewidth]{fig-5-3-bus-stress.pdf}
\caption{事件总线 publish 吞吐与丢帧率随订阅者数变化}
\label{fig:bus-stress}
\end{figure}

\textbf{零丢帧区间与广播 fan-out 物理本质}：1k/10k events/s 在任意 sub 数（最高 64）下均零丢帧，绝大多数 cell 的 p50 与 p99 处于亚微秒量级。这是 Tokio broadcast 内存 channel 在小负载下的典型表现：发布者将事件 \texttt{Arc::clone} 后写入环形缓冲，订阅者各自从缓冲读取——每次 fan-out 不复制数据本体只递增引用计数，因此「订阅者数翻倍」并不引起「fan-out 时间翻倍」，与 Tokio broadcast 文档「O(1) fan-out per subscriber」描述\cite{tokio2024broadcast}一致。

\textbf{唯一拐点：64 sub × 100k 事件/秒}。该 cell 触发约 42\% drops（13.6 M / (5 × 100 k × 64) ≈ 42.5\%），broadcast buffer（32k slots）不足以容纳 64 个 sub 的合计消费滞后；已接收事件的 p50 升至 163.5 us、p99 至 278.4 us。该现象为 SLA 公式中 $\max_i T_{consume,i}$ 项主导整体延迟下界提供了直接证据。Fuxi 选择 32768 slots 是工程经验折中：100k events/s 速率下对应 326 ms 缓冲窗口足以让 8 worker 个数级订阅者吸收偶发抖动；32k × 256 B = 8 MB 内存占用对现代机器可忽略。本实验中 64 sub × 100k 之所以丢帧是订阅者数量已超出工程经验上限——实际部署中订阅者通常等于 Worker 数加少量观察者，远低于 64。如需支持极端订阅者数可考虑由 broadcast 升级为分级 fan-out（broker 把订阅者按拓扑分组），思路与 ZooKeeper watcher 通知机制类似\cite{hunt2010zookeeper}。

\textbf{payload size 影响}：在零丢帧区间 small (256 B) 与 large (4096 B) 的 p50/p99 差异不超过 0.1 us——证实 Arc clone 而非 byte copy 的 fan-out 模型（byte copy 模型下 large 应比 small 慢约 16 倍）。\textbf{sustainable 容量与实际负载余量}：取 16 sub × 100k/s = 1.6 M events/s 处理量为 Fuxi 事件总线 sustainable 边界；Fuxi 实际工作负载（任务事件、工具调用、心跳、抄送）按 8 门客 × 单回合典型 50 事件 / 任务粒度约 1 s 估算合计约 100--1000 events/s，距 sustainable 边界还有 3 个数量级以上余量。这一余量保证系统负载高峰期事件总线不会成为瓶颈，且为未来突发流量场景（cc 门客在执行编辑工具时密集事件单回合可瞬间达数百 events/s、未来多用户并发会话或分布式跨节点合流\cite{corbett2012spanner}）留出空间。

\subsection{功能验收与本章小结}

除性能指标外本文还对系统关键交互路径进行功能验收，覆盖三类场景：busy 状态下用户消息不丢（玄女或门客 busy 时用户继续发送消息应入 pending queue 并在下一轮 drain 按序送达）、Codex 门客能起能派（codex 适配器能从 idle 启动、接受任务派发并写回结果，对应第三章关于多门客异构兼容的设计目标）、玄女不再通过 \texttt{fuxi status} 反复轮询（玄女通过订阅事件总线感知系统状态，对应公理「真实时不轮询」）。用户验收阶段还覆盖了抄送通路、门客派活后任务树迁移、对话区滚动、Resume 横幅等交互细节。系统在迭代过程中还修复了 cc 反连超时（Clash TUN 拦截 127.0.0.1，注入 NO\_PROXY）、TUI 启动期 stderr 撞 alt screen、idle GC 误杀玄女、cc/codex model fallback 硬编等若干典型工程缺陷，每一项都构成系统稳定性的边界条件。

\textbf{本章小结。} 本章围绕测试覆盖、吞吐量、延迟分布、扩展性、参数消融、事件总线压测、端到端延迟分解与功能验收七个维度对 Fuxi 进行了多组实测。定量结果可总结为五个数字（口径：合成基准衡量编排与通讯层在剥离 LLM 推理后的纯通讯吞吐上限，真实 LLM Agent 任务端到端吞吐受单次推理 1--30 s 主导）：8 worker × 10 ms 合成任务下达到 \textbf{665.00 tasks/s} 调度路径吞吐，工程损耗稳定在 17\% 量级；Worker 数 1$\to$16 范围内扩展效率始终保持在 \textbf{80.5\%--83.1\%}；任务派发 p50 为 \textbf{32.93 ms}、p99 为 36.66 ms，进程内事件总线广播 p50 仅 \textbf{0.08 ms}、p99 为 0.26 ms，两者相差三个数量级；事件总线 sustainable 容量为 \textbf{1.6 M events/s}，超过 Fuxi 实际工作负载三个数量级；端到端延迟在 1000 ms 任务粒度下通讯开销占比仅 \textbf{3.2\%}，工程上可忽略。这些数据共同验证了第二章关于「事件驱动通讯底座可同时兼顾吞吐、实时性与可追溯性」的设计假设。
